% GENERATED FILE. Edit canonical lessons, then run npm run generate:books.
% canonical-source-hash: fnv1a64:e8c2304081826460
% canonical-lessons: TE-C240-ettu, TE-C240-moyu, TE-C240-lagu, TE-C240-nettu, TE-C240-visuru

\chapter{To lift, To carry, To pull, To push, To throw}
\label{ch:te-to-lift-to-carry-to-pull-to-push-to-throw}
\hlchaptermodality{240}

\begin{chapteropening}
\textbf{By the end of this chapter:} \emph{I can say to lift, to carry, to pull, to push and to throw.}

Complete the last lesson of chapter 240: I can say to lift, to carry, to pull, to push and to throw.
\end{chapteropening}

\section[\texorpdfstring{\te{ఎత్తు} (ettu)}{ఎత్తు (ettu)}]{\te{ఎత్తు} (ettu) — to lift}

\label{lesson:TE-C240-ettu}



\begin{quote}
Before the new one: say the Telugu for a pension, then the Telugu for an agreement.
\end{quote}

\subsection*{You'll want to know: \te{ఎత్తు}}
\textbf{\te{ఎత్తు}} — \emph{ettu} — \textquotedblleft{}to lift\textquotedblright{}.

\begin{cousinweb}
\textbf{\te{ఎత్తు}} (\emph{ettu}) is native Dravidian. DEDR 796 lists \emph{ettu}, \textquotedblleft{}to lift up, raise\textquotedblright{}, with Kannada \emph{ettu}, \textquotedblleft{}to lift\textquotedblright{}.
\end{cousinweb}

\begin{grammarlens}[title={What you've built}]
One more word for this chapter.
\end{grammarlens}

\subsection*{Your turn}
\begin{itemize}
\raggedright
  \item \emph{Say it:} \emph{ettu}
  \item \emph{Say it:} \emph{ettu}, once more
  \item \emph{Say it:} say \emph{ettu}
  \item \emph{Recall:} say the Telugu for a pension, then the Telugu for an agreement, then say \emph{ettu} again
\end{itemize}

\begin{tcolorbox}[breakable,colback=teal!4,colframe=teal!35!black,title={Before you move on}]
What does \te{ఎత్తు} mean? (\textquotedblleft{}To lift\textquotedblright{}.) Say it once more.
\end{tcolorbox}
\section[\texorpdfstring{\te{మోయు} (mōyu)}{మోయు (mōyu)}]{\te{మోయు} (mōyu) — to carry (a load)}

\label{lesson:TE-C240-moyu}



\begin{quote}
Before the new one: say the Telugu for an agreement, then the Telugu for to lift.
\end{quote}

\subsection*{You'll want to know: \te{మోయు}}
\textbf{\te{మోయు}} — \emph{mōyu} — \textquotedblleft{}to carry (a load)\textquotedblright{}.

\begin{cousinweb}
\textbf{\te{మోయు}} (\emph{mōyu}) is native Dravidian. DEDR 5126 lists \emph{mōyu}, \textquotedblleft{}to bear, carry\textquotedblright{}, with Kannada \emph{mōṟ}, \textquotedblleft{}to carry on the head\textquotedblright{}. DEDR 5139 \emph{mōyu}, \textquotedblleft{}to be widowed\textquotedblright{}, is a different word.
\end{cousinweb}

\begin{grammarlens}[title={What you've built}]
One more word for this chapter.
\end{grammarlens}

\subsection*{Your turn}
\begin{itemize}
\raggedright
  \item \emph{Say it:} \emph{mōyu}
  \item \emph{Say it:} \emph{mōyu}, once more
  \item \emph{Say it:} say \emph{mōyu}
  \item \emph{Recall:} say the Telugu for an agreement, then the Telugu for to lift, then say \emph{mōyu} again
\end{itemize}

\begin{tcolorbox}[breakable,colback=teal!4,colframe=teal!35!black,title={Before you move on}]
What does \te{మోయు} mean? (\textquotedblleft{}To carry (a load)\textquotedblright{}.) Say it once more.
\end{tcolorbox}
\section[\texorpdfstring{\te{లాగు} (lāgu)}{లాగు (lāgu)}]{\te{లాగు} (lāgu) — to pull}

\label{lesson:TE-C240-lagu}



\begin{quote}
Before the new one: say the Telugu for to lift, then the Telugu for to carry.
\end{quote}

\subsection*{You'll want to know: \te{లాగు}}
\textbf{\te{లాగు}} — \emph{lāgu} — \textquotedblleft{}to pull\textquotedblright{}.

\begin{cousinweb}
The origin of \textbf{\te{లాగు}} (\emph{lāgu}) is uncertain. DEDR and CDIAL list no word of this shape meaning \textquotedblleft{}to pull\textquotedblright{}.
\end{cousinweb}

\begin{grammarlens}[title={What you've built}]
One more word for this chapter.
\end{grammarlens}

\subsection*{Your turn}
\begin{itemize}
\raggedright
  \item \emph{Say it:} \emph{lāgu}
  \item \emph{Say it:} \emph{lāgu}, once more
  \item \emph{Say it:} say \emph{lāgu}
  \item \emph{Recall:} say the Telugu for to lift, then the Telugu for to carry, then say \emph{lāgu} again
\end{itemize}

\begin{tcolorbox}[breakable,colback=teal!4,colframe=teal!35!black,title={Before you move on}]
What does \te{లాగు} mean? (\textquotedblleft{}To pull\textquotedblright{}.) Say it once more.
\end{tcolorbox}
\section[\texorpdfstring{\te{నెట్టు} (neṭṭu)}{నెట్టు (neṭṭu)}]{\te{నెట్టు} (neṭṭu) — to push}

\label{lesson:TE-C240-nettu}



\begin{quote}
Before the new one: say the Telugu for to carry, then the Telugu for to pull.
\end{quote}

\subsection*{You'll want to know: \te{నెట్టు}}
\textbf{\te{నెట్టు}} — \emph{neṭṭu} — \textquotedblleft{}to push\textquotedblright{}.

\begin{cousinweb}
\textbf{\te{నెట్టు}} (\emph{neṭṭu}) is native Dravidian. DEDR 3760 lists \emph{neṭṭu}, \textquotedblleft{}to push, shove\textquotedblright{}, with Tamil \emph{neṭṭu}, \textquotedblleft{}to thrust, push\textquotedblright{}.
\end{cousinweb}

\begin{grammarlens}[title={What you've built}]
One more word for this chapter.
\end{grammarlens}

\subsection*{Your turn}
\begin{itemize}
\raggedright
  \item \emph{Say it:} \emph{neṭṭu}
  \item \emph{Say it:} \emph{neṭṭu}, once more
  \item \emph{Say it:} say \emph{neṭṭu}
  \item \emph{Recall:} say the Telugu for to carry, then the Telugu for to pull, then say \emph{neṭṭu} again
\end{itemize}

\begin{tcolorbox}[breakable,colback=teal!4,colframe=teal!35!black,title={Before you move on}]
What does \te{నెట్టు} mean? (\textquotedblleft{}To push\textquotedblright{}.) Say it once more.
\end{tcolorbox}
\section[\texorpdfstring{\te{విసురు} (visuru)}{విసురు (visuru)}]{\te{విసురు} (visuru) — to throw, to fling}

\label{lesson:TE-C240-visuru}



\begin{quote}
Before the new one: say the Telugu for to pull, then the Telugu for to push.
\end{quote}

\subsection*{You'll want to know: \te{విసురు}}
\textbf{\te{విసురు}} — \emph{visuru} — \textquotedblleft{}to throw, to fling\textquotedblright{}.

\begin{cousinweb}
\textbf{\te{విసురు}} (\emph{visuru}) is native Dravidian. DEDR 5450 lists \emph{visuru}, \textquotedblleft{}to throw, fling\textquotedblright{}, with Tamil \emph{vīcu}, \textquotedblleft{}to throw, fling\textquotedblright{}. DEDR compares Sanskrit \emph{vyajana}, \textquotedblleft{}fanning; fan\textquotedblright{} (CDIAL 12043).
\end{cousinweb}

\begin{grammarlens}[title={What you've built}]
One more word for this chapter.
\end{grammarlens}

\subsection*{Your turn}
\begin{itemize}
\raggedright
  \item \emph{Say it:} \emph{visuru}
  \item \emph{Say it:} \emph{visuru}, once more
  \item \emph{Say it:} say \emph{visuru}
  \item \emph{Recall:} say the Telugu for to pull, then the Telugu for to push, then say \emph{visuru} again
\end{itemize}

\begin{tcolorbox}[breakable,colback=teal!4,colframe=teal!35!black,title={Before you move on}]
What does \te{విసురు} mean? (\textquotedblleft{}To throw, to fling\textquotedblright{}.) Say it once more.
\end{tcolorbox}
